\documentclass[12pt,a4paper]{article}

% ------------------------------------------------------------
% Кодировка, язык и шрифты
% ------------------------------------------------------------
\usepackage[T2A]{fontenc}
\usepackage[utf8]{inputenc}
\usepackage[russian]{babel}

% ------------------------------------------------------------
% Математика
% ------------------------------------------------------------
\usepackage{amsmath,amssymb,amsthm,mathtools}
\usepackage{bm}

% ------------------------------------------------------------
% Оформление
% ------------------------------------------------------------
\usepackage[a4paper,margin=2cm]{geometry}
\usepackage[expansion=false]{microtype}
\usepackage{booktabs}
\usepackage{array}
\usepackage{longtable}
\usepackage{enumitem}
\usepackage{algorithm}
\usepackage{algpseudocode}
\usepackage{tikz}
\usetikzlibrary{positioning,arrows.meta,calc}
\usepackage[hidelinks]{hyperref}

% ------------------------------------------------------------
% Теоремы и определения
% ------------------------------------------------------------
\newtheorem{theorem}{Теорема}[section]
\newtheorem{proposition}[theorem]{Предложение}
\newtheorem{lemma}[theorem]{Лемма}

\theoremstyle{definition}
\newtheorem{definition}[theorem]{Определение}
\newtheorem{example}[theorem]{Пример}
\newtheorem{remark}[theorem]{Замечание}

% ------------------------------------------------------------
% Обозначения
% ------------------------------------------------------------
\newcommand{\R}{\mathbb{R}}
\newcommand{\T}{\mathsf{T}}
\newcommand{\E}{\mathbb{E}}
\newcommand{\1}{\mathbf{1}}
\newcommand{\eps}{\varepsilon}
\newcommand{\dd}{\,\mathrm{d}}

\DeclareMathOperator{\sign}{sign}
\DeclareMathOperator{\ReLU}{ReLU}
\DeclareMathOperator{\softmax}{softmax}
\DeclareMathOperator{\argmin}{arg\,min}
\DeclareMathOperator{\argmax}{arg\,max}
\DeclareMathOperator{\dist}{dist}

% ------------------------------------------------------------
% Локализация алгоритмов
% ------------------------------------------------------------
\floatname{algorithm}{Алгоритм}
\algrenewcommand\algorithmicrequire{\textbf{Вход:}}
\algrenewcommand\algorithmicensure{\textbf{Выход:}}
\algrenewcommand\algorithmicwhile{\textbf{пока}}
\algrenewcommand\algorithmicdo{\textbf{выполнять}}
\algrenewcommand\algorithmicfor{\textbf{для}}
\algrenewcommand\algorithmicforall{\textbf{для всех}}
\algrenewcommand\algorithmicif{\textbf{если}}
\algrenewcommand\algorithmicthen{\textbf{то}}
\algrenewcommand\algorithmicelse{\textbf{иначе}}
\algrenewcommand\algorithmicend{\textbf{конец}}
\algrenewcommand\algorithmicreturn{\textbf{вернуть}}

\title{Билет №96 \\ \small Нейронные сети. Метод опорных векторов. (ИТМО)}
\author{Сериков Владислав Александрович}
\date{19 июля 2026}

\begin{document}

\maketitle
\tableofcontents
\newpage

% ============================================================
\section{Общая постановка задачи машинного обучения}
% ============================================================

Пусть задана обучающая выборка
\[
    \mathcal{D}
    =
    \{(\bm{x}_i,y_i)\}_{i=1}^{n},
    \qquad
    \bm{x}_i\in\mathcal{X},
    \quad
    y_i\in\mathcal{Y}.
\]
Вектор $\bm{x}_i$ содержит признаки объекта, а $y_i$ является целевой
переменной.

В зависимости от вида множества $\mathcal{Y}$ различают:

\begin{itemize}
    \item \textbf{регрессию}: $\mathcal{Y}\subseteq\R$;
    \item \textbf{бинарную классификацию}: $\mathcal{Y}=\{-1,+1\}$
    или $\mathcal{Y}=\{0,1\}$;
    \item \textbf{многоклассовую классификацию}:
    $\mathcal{Y}=\{1,\ldots,K\}$.
\end{itemize}

Модель $f(\bm{x};\bm{\theta})$ зависит от обучаемых параметров
$\bm{\theta}$. Качество предсказания измеряется функцией потерь
\[
    \ell\bigl(y,f(\bm{x};\bm{\theta})\bigr).
\]

\begin{definition}
\textbf{Эмпирическим риском} называется среднее значение функции потерь
на обучающей выборке:
\[
    \widehat{R}(\bm{\theta})
    =
    \frac{1}{n}
    \sum_{i=1}^{n}
    \ell\bigl(y_i,f(\bm{x}_i;\bm{\theta})\bigr).
\]
\end{definition}

Обучение обычно сводится к решению задачи
\[
    \bm{\theta}^{*}
    =
    \argmin_{\bm{\theta}}
    \left[
        \widehat{R}(\bm{\theta})
        +
        \lambda \Omega(\bm{\theta})
    \right],
\]
где $\Omega$ --- регуляризатор, ограничивающий сложность модели,
а $\lambda\geq 0$ --- коэффициент регуляризации.

Важно различать:

\begin{itemize}
    \item \textbf{параметры} --- изменяются алгоритмом обучения;
    \item \textbf{гиперпараметры} --- задаются до обучения или выбираются
    по валидационной выборке;
    \item \textbf{обобщающую способность} --- качество модели на новых
    объектах, не участвовавших в обучении.
\end{itemize}

Обычно данные разделяют на обучающую, валидационную и тестовую части.
Тестовая выборка не должна использоваться для настройки модели.

% ============================================================
\section{Искусственный нейрон и перцептрон}
% ============================================================

\subsection{Математическая модель нейрона}

Искусственный нейрон вычисляет взвешенную сумму входов:
\[
    z
    =
    \sum_{j=1}^{d}w_jx_j+b
    =
    \bm{w}^{\T}\bm{x}+b,
\]
после чего применяет функцию активации:
\[
    a=\varphi(z).
\]

Здесь $\bm{w}\in\R^d$ --- вектор весов, $b\in\R$ --- смещение,
$\varphi$ --- функция активации.

\begin{figure}[h]
\centering
\begin{tikzpicture}[
    input/.style={circle,draw,minimum size=8mm},
    neuron/.style={circle,draw,minimum size=13mm,fill=gray!10},
    >=Latex
]
    \node[input] (x1) at (0,1.5) {$x_1$};
    \node[input] (x2) at (0,0.5) {$x_2$};
    \node (dots) at (0,-0.3) {$\vdots$};
    \node[input] (xd) at (0,-1.3) {$x_d$};

    \node[neuron] (n) at (3,0) {$\varphi$};
    \node (out) at (5,0) {$a$};

    \draw[->] (x1) -- node[above,sloped] {$w_1$} (n);
    \draw[->] (x2) -- node[above,sloped] {$w_2$} (n);
    \draw[->] (xd) -- node[below,sloped] {$w_d$} (n);
    \draw[->] (n) -- (out);
    \node[below=3mm of n] {$z=\bm{w}^{\T}\bm{x}+b$};
\end{tikzpicture}
\caption{Математическая модель искусственного нейрона.}
\end{figure}

\subsection{Основные функции активации}

\begin{center}
\begin{longtable}{p{2.7cm} p{5.2cm} p{5.2cm}}
\toprule
\textbf{Функция} & \textbf{Формула} & \textbf{Производная} \\
\midrule
Пороговая
&
\[
\varphi(z)=
\begin{cases}
1,&z\geq 0,\\
0,&z<0
\end{cases}
\]
&
Не дифференцируема в нуле, вне нуля производная равна $0$.
\\
\midrule
Сигмоида
&
\[
\sigma(z)=\frac{1}{1+e^{-z}}
\]
&
\[
\sigma'(z)=\sigma(z)\bigl(1-\sigma(z)\bigr)
\]
\\
\midrule
Гиперболический тангенс
&
\[
\tanh z=\frac{e^z-e^{-z}}{e^z+e^{-z}}
\]
&
\[
\frac{\dd}{\dd z}\tanh z=1-\tanh^2z
\]
\\
\midrule
ReLU
&
\[
\ReLU(z)=\max(0,z)
\]
&
\[
\ReLU'(z)=
\begin{cases}
0,&z<0,\\
1,&z>0.
\end{cases}
\]
В точке $z=0$ обычно выбирают одно из субградиентных значений.
\\
\midrule
Leaky ReLU
&
\[
\varphi(z)=
\begin{cases}
z,&z\geq 0,\\
\alpha z,&z<0
\end{cases}
\]
&
\[
\varphi'(z)=
\begin{cases}
1,&z>0,\\
\alpha,&z<0.
\end{cases}
\]
\\
\bottomrule
\end{longtable}
\end{center}

Без нелинейной функции активации композиция нескольких линейных слоёв
остаётся линейным преобразованием:
\[
    W_2(W_1\bm{x}+\bm{b}_1)+\bm{b}_2
    =
    (W_2W_1)\bm{x}
    +
    (W_2\bm{b}_1+\bm{b}_2).
\]
Поэтому нелинейность является принципиально важной составляющей
многослойной нейронной сети.

\subsection{Линейный классификатор}

В бинарной классификации линейный классификатор имеет вид
\[
    \widehat{y}
    =
    \sign(\bm{w}^{\T}\bm{x}+b).
\]

Гиперплоскость
\[
    \bm{w}^{\T}\bm{x}+b=0
\]
разделяет пространство признаков на два полупространства.

\begin{definition}
Выборка $\{(\bm{x}_i,y_i)\}_{i=1}^{n}$, где $y_i\in\{-1,+1\}$,
называется \textbf{линейно разделимой}, если существуют
$\bm{w}$ и $b$, для которых
\[
    y_i(\bm{w}^{\T}\bm{x}_i+b)>0
    \qquad
    \text{для всех }i.
\]
\end{definition}

\subsection{Алгоритм обучения перцептрона}

Перцептрон изменяет параметры только на ошибочно классифицированных
объектах.

\begin{algorithm}[h]
\caption{Обучение перцептрона}
\begin{algorithmic}[1]
\Require Выборка
$\{(\bm{x}_i,y_i)\}_{i=1}^{n}$,
$y_i\in\{-1,+1\}$,
скорость обучения $\eta>0$.
\Ensure Вектор весов $\bm{w}$ и смещение $b$.
\State Инициализировать $\bm{w}\gets\bm{0}$, $b\gets 0$.
\Repeat
    \State $\textit{errors}\gets 0$.
    \For{$i=1,\ldots,n$}
        \If{$y_i(\bm{w}^{\T}\bm{x}_i+b)\leq 0$}
            \State $\bm{w}\gets\bm{w}+\eta y_i\bm{x}_i$.
            \State $b\gets b+\eta y_i$.
            \State $\textit{errors}\gets\textit{errors}+1$.
        \EndIf
    \EndFor
\Until{$\textit{errors}=0$ или достигнуто ограничение числа эпох}
\State \Return $(\bm{w},b)$.
\end{algorithmic}
\end{algorithm}

Вектор смещения можно включить в вектор весов:
\[
    \widetilde{\bm{x}}
    =
    \begin{pmatrix}
        \bm{x}\\
        1
    \end{pmatrix},
    \qquad
    \widetilde{\bm{w}}
    =
    \begin{pmatrix}
        \bm{w}\\
        b
    \end{pmatrix}.
\]
Тогда
\[
    \widetilde{\bm{w}}^{\T}\widetilde{\bm{x}}
    =
    \bm{w}^{\T}\bm{x}+b.
\]

\subsection{Теорема о сходимости перцептрона}

\begin{theorem}[Сходимость перцептрона]
Пусть дана выборка
\[
    \{(\bm{x}_i,y_i)\}_{i=1}^{n},
    \qquad
    y_i\in\{-1,+1\},
\]
и существуют единичный вектор $\bm{u}$ и число $\gamma>0$ такие, что
\[
    y_i\bm{u}^{\T}\bm{x}_i\geq\gamma
    \qquad
    \text{для всех }i.
\]
Пусть также
\[
    \|\bm{x}_i\|\leq R
    \qquad
    \text{для всех }i.
\]
Если перцептрон начинает с $\bm{w}_0=\bm{0}$ и при каждой ошибке
выполняет обновление
\[
    \bm{w}_{t+1}
    =
    \bm{w}_t+y_i\bm{x}_i,
\]
то общее число ошибок $M$ удовлетворяет оценке
\[
    M\leq\frac{R^2}{\gamma^2}.
\]
В частности, на линейно разделимой выборке алгоритм совершает конечное
число ошибок.
\end{theorem}

\begin{proof}
Рассмотрим только моменты, в которые перцептрон совершает ошибку.

После очередного ошибочного предсказания:
\[
    \bm{w}_{t+1}
    =
    \bm{w}_t+y_i\bm{x}_i.
\]
Умножим это равенство скалярно на разделяющий единичный вектор
$\bm{u}$:
\[
    \bm{u}^{\T}\bm{w}_{t+1}
    =
    \bm{u}^{\T}\bm{w}_t
    +
    y_i\bm{u}^{\T}\bm{x}_i.
\]
По условию разделимости
\[
    y_i\bm{u}^{\T}\bm{x}_i\geq\gamma.
\]
Следовательно,
\[
    \bm{u}^{\T}\bm{w}_{t+1}
    \geq
    \bm{u}^{\T}\bm{w}_t+\gamma.
\]
Так как $\bm{w}_0=\bm{0}$, после $M$ ошибок получаем
\[
    \bm{u}^{\T}\bm{w}_M\geq M\gamma.
\]

Теперь оценим норму $\bm{w}_M$. При ошибке выполняется
\[
    y_i\bm{w}_t^{\T}\bm{x}_i\leq 0.
\]
Поэтому
\begin{align*}
    \|\bm{w}_{t+1}\|^2
    &=
    \|\bm{w}_t+y_i\bm{x}_i\|^2\\
    &=
    \|\bm{w}_t\|^2
    +
    2y_i\bm{w}_t^{\T}\bm{x}_i
    +
    y_i^2\|\bm{x}_i\|^2\\
    &\leq
    \|\bm{w}_t\|^2+R^2,
\end{align*}
поскольку $y_i^2=1$. После $M$ ошибок:
\[
    \|\bm{w}_M\|^2\leq MR^2,
    \qquad
    \|\bm{w}_M\|\leq R\sqrt{M}.
\]

По неравенству Коши--Буняковского и условию $\|\bm{u}\|=1$:
\[
    \bm{u}^{\T}\bm{w}_M
    \leq
    \|\bm{u}\|\,\|\bm{w}_M\|
    \leq
    R\sqrt{M}.
\]
Объединяя две оценки, получаем
\[
    M\gamma
    \leq
    R\sqrt{M}.
\]
Если $M>0$, то после деления на $\sqrt{M}$:
\[
    \sqrt{M}\,\gamma\leq R.
\]
Следовательно,
\[
    M\leq\frac{R^2}{\gamma^2}.
\]
\end{proof}

\subsection{Примеры линейной разделимости}

\begin{example}[Логическая функция AND]
Пусть $\bm{x}=(x_1,x_2)$, где $x_1,x_2\in\{0,1\}$, а положительным
является только объект $(1,1)$. Классификатор
\[
    \widehat{y}
    =
    \sign\left(x_1+x_2-\frac{3}{2}\right)
\]
правильно реализует функцию AND.
\end{example}

\begin{example}[Неразделимость XOR]
Для XOR отрицательными являются точки $(0,0)$ и $(1,1)$,
а положительными --- $(1,0)$ и $(0,1)$.

Предположим, что существует линейный разделитель
\[
    f(x_1,x_2)=w_1x_1+w_2x_2+b.
\]
Из отрицательности точки $(0,0)$ следует $b<0$. Из положительности
точек $(1,0)$ и $(0,1)$:
\[
    w_1+b>0,
    \qquad
    w_2+b>0.
\]
Значит,
\[
    w_1>-b,
    \qquad
    w_2>-b.
\]
Тогда
\[
    w_1+w_2+b>-b>0,
\]
но точка $(1,1)$ должна быть отрицательной. Получено противоречие.
Следовательно, XOR не является линейно разделимой задачей.
\end{example}

% ============================================================
\section{Многослойные нейронные сети}
% ============================================================

\subsection{Полносвязная сеть}

Многослойный перцептрон представляет собой композицию функций.
Положим
\[
    \bm{a}^{(0)}=\bm{x}.
\]
Для слоя $\ell=1,\ldots,L$ вычисляются
\[
    \bm{z}^{(\ell)}
    =
    W^{(\ell)}\bm{a}^{(\ell-1)}
    +
    \bm{b}^{(\ell)},
\]
\[
    \bm{a}^{(\ell)}
    =
    \varphi^{(\ell)}\bigl(\bm{z}^{(\ell)}\bigr).
\]
Функция активации применяется покомпонентно.

Если слой $\ell-1$ содержит $n_{\ell-1}$ нейронов, а слой $\ell$ ---
$n_\ell$ нейронов, то
\[
    W^{(\ell)}
    \in
    \R^{n_\ell\times n_{\ell-1}},
    \qquad
    \bm{b}^{(\ell)}
    \in
    \R^{n_\ell}.
\]
Число параметров слоя равно
\[
    n_\ell n_{\ell-1}+n_\ell.
\]

\begin{figure}[h]
\centering
\begin{tikzpicture}[
    neuron/.style={circle,draw,minimum size=8mm},
    >=Latex
]
    \foreach \i/\y in {1/1.3,2/0,3/-1.3}
        \node[neuron] (x\i) at (0,\y) {$x_{\i}$};

    \foreach \i/\y in {1/1.8,2/0.6,3/-0.6,4/-1.8}
        \node[neuron,fill=blue!8] (h\i) at (3,\y) {$h_{\i}$};

    \foreach \i/\y in {1/0.7,2/-0.7}
        \node[neuron,fill=green!8] (o\i) at (6,\y) {$o_{\i}$};

    \foreach \i in {1,2,3}
        \foreach \j in {1,2,3,4}
            \draw[->,gray] (x\i) -- (h\j);

    \foreach \i in {1,2,3,4}
        \foreach \j in {1,2}
            \draw[->,gray] (h\i) -- (o\j);

    \node at (0,-2.5) {входной слой};
    \node at (3,-2.5) {скрытый слой};
    \node at (6,-2.5) {выходной слой};
\end{tikzpicture}
\caption{Полносвязная нейронная сеть.}
\end{figure}

\subsection{Реализация XOR сетью с ReLU}

Однослойный линейный классификатор не реализует XOR, но это делает
сеть с одним скрытым слоем:
\[
    h_1=\ReLU(x_1+x_2),
    \qquad
    h_2=\ReLU(x_1+x_2-1),
\]
\[
    N(x_1,x_2)=h_1-2h_2.
\]

\begin{center}
\begin{tabular}{ccccc}
\toprule
$x_1$ & $x_2$ & $h_1$ & $h_2$ & $N(x_1,x_2)$\\
\midrule
0 & 0 & 0 & 0 & 0\\
0 & 1 & 1 & 0 & 1\\
1 & 0 & 1 & 0 & 1\\
1 & 1 & 2 & 1 & 0\\
\bottomrule
\end{tabular}
\end{center}

Таким образом, нелинейный скрытый слой строит новое представление
объектов, в котором задача становится разделимой.

\subsection{Универсальная аппроксимация на отрезке}

Следующий результат показывает, что даже сеть с одним скрытым слоем
ReLU может приближать произвольные непрерывные функции одной
переменной.

\begin{theorem}[Универсальная аппроксимация непрерывных функций на отрезке]
Пусть $f\in C([a,b])$. Тогда для любого $\eps>0$ существует функция
вида
\[
    N(x)
    =
    c_0+
    \sum_{k=0}^{m}c_{k+1}\ReLU(x-t_k),
\]
для которой
\[
    \sup_{x\in[a,b]}|f(x)-N(x)|<\eps.
\]
Следовательно, функции, вычисляемые сетями с одним скрытым слоем ReLU,
плотны в $C([a,b])$ в равномерной норме.
\end{theorem}

\begin{proof}
Поскольку функция $f$ непрерывна на компакте $[a,b]$, она равномерно
непрерывна. Поэтому для заданного $\eps>0$ существует $\delta>0$ такое,
что
\[
    |x-y|<\delta
    \quad\Longrightarrow\quad
    |f(x)-f(y)|<\eps.
\]

Выберем разбиение
\[
    a=t_0<t_1<\dots<t_m=b
\]
с диаметром
\[
    \max_{0\leq k<m}(t_{k+1}-t_k)<\delta.
\]
Построим кусочно-линейный интерполянт $p$, удовлетворяющий условиям
\[
    p(t_k)=f(t_k).
\]

На отрезке $[t_k,t_{k+1}]$ положим
\[
    \lambda
    =
    \frac{x-t_k}{t_{k+1}-t_k}.
\]
Тогда $0\leq\lambda\leq1$ и
\[
    p(x)
    =
    (1-\lambda)f(t_k)+\lambda f(t_{k+1}).
\]
Следовательно,
\begin{align*}
    |f(x)-p(x)|
    &=
    \left|
        (1-\lambda)(f(x)-f(t_k))
        +
        \lambda(f(x)-f(t_{k+1}))
    \right|\\
    &\leq
    (1-\lambda)|f(x)-f(t_k)|
    +
    \lambda|f(x)-f(t_{k+1})|.
\end{align*}
Так как
\[
    |x-t_k|<\delta,
    \qquad
    |x-t_{k+1}|<\delta,
\]
то оба слагаемых меньше $\eps$, а значит,
\[
    |f(x)-p(x)|<\eps.
\]
Таким образом,
\[
    \sup_{x\in[a,b]}|f(x)-p(x)|<\eps.
\]

Осталось показать, что любой такой кусочно-линейный интерполянт
выражается через ReLU. Обозначим наклон на отрезке
$[t_k,t_{k+1}]$ через
\[
    s_k
    =
    \frac{f(t_{k+1})-f(t_k)}{t_{k+1}-t_k}.
\]
Рассмотрим функцию
\[
    N(x)
    =
    f(a)
    +
    s_0\ReLU(x-a)
    +
    \sum_{k=1}^{m-1}
        (s_k-s_{k-1})\ReLU(x-t_k).
\]

При $x\in[t_j,t_{j+1}]$ её производная в точках, отличных от узлов,
равна
\begin{align*}
    N'(x)
    &=
    s_0+
    \sum_{k=1}^{j}(s_k-s_{k-1})\\
    &=
    s_j.
\end{align*}
Кроме того,
\[
    N(a)=f(a).
\]
Следовательно, $N$ имеет на каждом отрезке тот же наклон и то же
начальное значение, что и интерполянт $p$. Поэтому $N=p$ на всём
$[a,b]$.

Получаем
\[
    \sup_{x\in[a,b]}|f(x)-N(x)|<\eps.
\]
\end{proof}

\begin{remark}
Универсальная аппроксимация является утверждением о существовании
параметров. Она сама по себе не гарантирует, что параметры будут
быстро найдены градиентным методом, что сеть будет компактной или
что она хорошо обобщится на новые данные.
\end{remark}

% ============================================================
\section{Выходные слои и функции потерь}
% ============================================================

\subsection{Регрессия}

Для вещественного предсказания выходной слой обычно является линейным:
\[
    \widehat{y}=z.
\]
Часто используется квадратичная функция потерь:
\[
    \ell(y,\widehat{y})
    =
    \frac{1}{2}(y-\widehat{y})^2.
\]
Её производная:
\[
    \frac{\partial\ell}{\partial\widehat{y}}
    =
    \widehat{y}-y.
\]

Квадратичная потеря особенно чувствительна к выбросам. Более устойчива
абсолютная потеря
\[
    \ell(y,\widehat{y})=|y-\widehat{y}|,
\]
но она не дифференцируема в нуле.

\subsection{Бинарная классификация}

Пусть выходной логит равен
\[
    z=\bm{w}^{\T}\bm{h}+b.
\]
Вероятность положительного класса задаётся сигмоидой:
\[
    p=P(y=1\mid\bm{x})=\sigma(z).
\]

Бинарная кросс-энтропия:
\[
    \ell(y,p)
    =
    -y\log p-(1-y)\log(1-p).
\]
Для композиции сигмоиды и кросс-энтропии:
\[
    \frac{\partial\ell}{\partial z}=p-y.
\]

\subsection{Многоклассовая классификация}

Для $K$ классов сеть выдаёт вектор логитов
\[
    \bm{z}=(z_1,\ldots,z_K).
\]
Функция softmax преобразует его в распределение вероятностей:
\[
    p_k
    =
    \frac{e^{z_k}}
         {\sum_{j=1}^{K}e^{z_j}}.
\]
Для численной устойчивости используют эквивалентную формулу
\[
    p_k
    =
    \frac{e^{z_k-z_{\max}}}
         {\sum_{j=1}^{K}e^{z_j-z_{\max}}},
    \qquad
    z_{\max}=\max_jz_j.
\]

Если $\bm{y}$ --- one-hot-вектор истинного класса, то категориальная
кросс-энтропия имеет вид
\[
    \ell(\bm{y},\bm{p})
    =
    -\sum_{k=1}^{K}y_k\log p_k.
\]

Производная softmax равна
\[
    \frac{\partial p_k}{\partial z_j}
    =
    p_k(\delta_{kj}-p_j),
\]
где $\delta_{kj}$ --- символ Кронекера. Поэтому
\begin{align*}
    \frac{\partial\ell}{\partial z_j}
    &=
    -\sum_{k=1}^{K}
    \frac{y_k}{p_k}
    \frac{\partial p_k}{\partial z_j}\\
    &=
    -\sum_{k=1}^{K}y_k(\delta_{kj}-p_j)\\
    &=
    p_j-y_j,
\end{align*}
поскольку $\sum_k y_k=1$. В векторной форме:
\[
    \nabla_{\bm{z}}\ell
    =
    \bm{p}-\bm{y}.
\]

% ============================================================
\section{Обратное распространение ошибки}
% ============================================================

\subsection{Общая идея}

Обратное распространение ошибки, или backpropagation, представляет
собой эффективное применение правила дифференцирования сложной функции
к вычислительному графу нейронной сети.

Пусть
\[
    \bm{z}^{(\ell)}
    =
    W^{(\ell)}\bm{a}^{(\ell-1)}
    +
    \bm{b}^{(\ell)},
    \qquad
    \bm{a}^{(\ell)}
    =
    \varphi^{(\ell)}(\bm{z}^{(\ell)}).
\]
Обозначим
\[
    \bm{\delta}^{(\ell)}
    =
    \frac{\partial J}{\partial\bm{z}^{(\ell)}}.
\]

Для последнего слоя вектор $\bm{\delta}^{(L)}$ определяется сочетанием
выходной функции активации и функции потерь. Например, для softmax
с кросс-энтропией:
\[
    \bm{\delta}^{(L)}
    =
    \widehat{\bm{y}}-\bm{y}.
\]

Для скрытых слоёв:
\[
    \bm{\delta}^{(\ell)}
    =
    \left(W^{(\ell+1)}\right)^{\T}
    \bm{\delta}^{(\ell+1)}
    \odot
    \varphi'^{(\ell)}(\bm{z}^{(\ell)}),
\]
где $\odot$ обозначает покомпонентное умножение.

Градиенты параметров имеют вид
\[
    \frac{\partial J}{\partial W^{(\ell)}}
    =
    \bm{\delta}^{(\ell)}
    \left(\bm{a}^{(\ell-1)}\right)^{\T},
\]
\[
    \frac{\partial J}{\partial\bm{b}^{(\ell)}}
    =
    \bm{\delta}^{(\ell)}.
\]

\subsection{Вывод формул}

Для слоя
\[
    \bm{z}=W\bm{a}+\bm{b}
\]
дифференциал равен
\[
    \dd\bm{z}
    =
    \dd W\,\bm{a}
    +
    W\,\dd\bm{a}
    +
    \dd\bm{b}.
\]
Если
\[
    \bm{\delta}
    =
    \frac{\partial J}{\partial\bm{z}},
\]
то
\begin{align*}
    \dd J
    &=
    \bm{\delta}^{\T}\dd\bm{z}\\
    &=
    \bm{\delta}^{\T}\dd W\,\bm{a}
    +
    \bm{\delta}^{\T}W\,\dd\bm{a}
    +
    \bm{\delta}^{\T}\dd\bm{b}.
\end{align*}

Первое слагаемое можно записать через след:
\[
    \bm{\delta}^{\T}\dd W\,\bm{a}
    =
    \operatorname{tr}
    \left(
        (\bm{\delta}\bm{a}^{\T})^{\T}
        \dd W
    \right).
\]
Следовательно,
\[
    \frac{\partial J}{\partial W}
    =
    \bm{\delta}\bm{a}^{\T}.
\]
Из остальных слагаемых:
\[
    \frac{\partial J}{\partial\bm{a}}
    =
    W^{\T}\bm{\delta},
    \qquad
    \frac{\partial J}{\partial\bm{b}}
    =
    \bm{\delta}.
\]

Если
\[
    \bm{a}
    =
    \varphi(\bm{z}_{\mathrm{prev}}),
\]
то по правилу цепочки
\[
    \frac{\partial J}{\partial\bm{z}_{\mathrm{prev}}}
    =
    \frac{\partial J}{\partial\bm{a}}
    \odot
    \varphi'(\bm{z}_{\mathrm{prev}}),
\]
что и даёт рекуррентную формулу обратного распространения.

\subsection{Алгоритм обучения сети}

\begin{algorithm}[h]
\caption{Обучение нейронной сети методом обратного распространения}
\begin{algorithmic}[1]
\Require Выборка $\mathcal{D}$, архитектура сети, функция потерь,
число эпох $T$, размер мини-пакета $B$.
\Ensure Обученные параметры $\bm{\theta}$.
\State Инициализировать матрицы весов и смещения.
\For{$t=1,\ldots,T$}
    \State Случайно перемешать обучающую выборку.
    \For{каждого мини-пакета}
        \State Выполнить прямой проход и сохранить
        $\bm{z}^{(\ell)}$, $\bm{a}^{(\ell)}$.
        \State Вычислить значение функции потерь.
        \State Вычислить ошибку выходного слоя
        $\bm{\delta}^{(L)}$.
        \For{$\ell=L-1,L-2,\ldots,1$}
            \State
            $\bm{\delta}^{(\ell)}
            \gets
            (W^{(\ell+1)})^{\T}
            \bm{\delta}^{(\ell+1)}
            \odot
            \varphi'^{(\ell)}(\bm{z}^{(\ell)})$.
        \EndFor
        \For{$\ell=1,\ldots,L$}
            \State Вычислить
            $\nabla_{W^{(\ell)}}J$ и
            $\nabla_{\bm{b}^{(\ell)}}J$.
            \State Обновить параметры выбранным оптимизатором.
        \EndFor
    \EndFor
\EndFor
\State \Return $\bm{\theta}$.
\end{algorithmic}
\end{algorithm}

\subsection{Численный пример одного шага}

Рассмотрим сеть с одним входом и одним скрытым нейроном:
\[
    z_1=w_1x+b_1,
    \qquad
    h=\ReLU(z_1),
\]
\[
    \widehat{y}=w_2h+b_2,
    \qquad
    J=\frac{1}{2}(\widehat{y}-y)^2.
\]

Пусть
\[
    x=2,
    \quad
    y=1,
    \quad
    w_1=0.5,
    \quad
    b_1=0,
    \quad
    w_2=0.5,
    \quad
    b_2=0.
\]

Прямой проход:
\[
    z_1=0.5\cdot2=1,
    \qquad
    h=\ReLU(1)=1,
\]
\[
    \widehat{y}=0.5\cdot1=0.5,
\]
\[
    J=\frac{1}{2}(0.5-1)^2=0.125.
\]

Ошибка выходного слоя:
\[
    \delta_2
    =
    \frac{\partial J}{\partial\widehat{y}}
    =
    \widehat{y}-y=-0.5.
\]
Поэтому
\[
    \frac{\partial J}{\partial w_2}
    =
    \delta_2h=-0.5,
    \qquad
    \frac{\partial J}{\partial b_2}
    =
    \delta_2=-0.5.
\]

Так как $z_1>0$, имеем $\ReLU'(z_1)=1$:
\[
    \delta_1
    =
    \delta_2w_2\ReLU'(z_1)
    =
    -0.5\cdot0.5=-0.25.
\]
Следовательно,
\[
    \frac{\partial J}{\partial w_1}
    =
    \delta_1x=-0.5,
    \qquad
    \frac{\partial J}{\partial b_1}
    =
    -0.25.
\]

При скорости обучения $\eta=0.1$:
\[
    w_2^{\mathrm{new}}
    =
    0.5-0.1(-0.5)=0.55,
\]
\[
    b_2^{\mathrm{new}}=0.05,
    \qquad
    w_1^{\mathrm{new}}=0.55,
    \qquad
    b_1^{\mathrm{new}}=0.025.
\]

После обновления:
\[
    z_1^{\mathrm{new}}=0.55\cdot2+0.025=1.125,
\]
\[
    \widehat{y}^{\mathrm{new}}
    =
    0.55\cdot1.125+0.05
    =
    0.66875.
\]
Новое значение потерь:
\[
    J_{\mathrm{new}}
    =
    \frac{1}{2}(0.66875-1)^2
    \approx0.0549,
\]
то есть один градиентный шаг уменьшил ошибку.

% ============================================================
\section{Оптимизация нейронных сетей}
% ============================================================

\subsection{Градиентный спуск}

Полный градиентный спуск использует всю выборку:
\[
    \bm{\theta}_{t+1}
    =
    \bm{\theta}_t
    -
    \eta\nabla_{\bm{\theta}}
    \widehat{R}(\bm{\theta}_t).
\]

Стохастический градиентный спуск использует один случайный объект:
\[
    \bm{\theta}_{t+1}
    =
    \bm{\theta}_t
    -
    \eta
    \nabla_{\bm{\theta}}
    \ell_i(\bm{\theta}_t).
\]

На практике чаще используется мини-пакет $\mathcal{B}_t$:
\[
    \bm{g}_t
    =
    \frac{1}{|\mathcal{B}_t|}
    \sum_{i\in\mathcal{B}_t}
    \nabla_{\bm{\theta}}\ell_i(\bm{\theta}_t),
\]
\[
    \bm{\theta}_{t+1}
    =
    \bm{\theta}_t-\eta\bm{g}_t.
\]

Мини-пакеты обеспечивают компромисс между точностью оценки градиента,
вычислительной эффективностью и объёмом используемой памяти.

\subsection{Метод момента}

Метод момента накапливает экспоненциально сглаженное направление:
\[
    \bm{v}_t
    =
    \beta\bm{v}_{t-1}
    +
    \bm{g}_t,
\]
\[
    \bm{\theta}_{t+1}
    =
    \bm{\theta}_t-\eta\bm{v}_t.
\]
Параметр $\beta$ обычно выбирается близким к единице. Метод уменьшает
осцилляции и ускоряет движение в направлениях устойчивого градиента.

\subsection{Оптимизатор Adam}

Adam хранит оценки первого и второго моментов градиента.

\begin{algorithm}[h]
\caption{Adam}
\begin{algorithmic}[1]
\Require Начальные параметры $\bm{\theta}_0$, скорость обучения
$\eta$, коэффициенты $\beta_1,\beta_2\in[0,1)$,
число $\eps_{\mathrm{Adam}}>0$.
\Ensure Параметры $\bm{\theta}$.
\State $\bm{m}_0\gets\bm{0}$, $\bm{v}_0\gets\bm{0}$.
\For{$t=1,2,\ldots$}
    \State Вычислить градиент
    $\bm{g}_t\gets\nabla_{\bm{\theta}}J_t(\bm{\theta}_{t-1})$.
    \State
    $\bm{m}_t
    \gets
    \beta_1\bm{m}_{t-1}
    +(1-\beta_1)\bm{g}_t$.
    \State
    $\bm{v}_t
    \gets
    \beta_2\bm{v}_{t-1}
    +(1-\beta_2)(\bm{g}_t\odot\bm{g}_t)$.
    \State
    $\widehat{\bm{m}}_t
    \gets
    \bm{m}_t/(1-\beta_1^t)$.
    \State
    $\widehat{\bm{v}}_t
    \gets
    \bm{v}_t/(1-\beta_2^t)$.
    \State
    $\bm{\theta}_t
    \gets
    \bm{\theta}_{t-1}
    -
    \eta
    \widehat{\bm{m}}_t/
    (\sqrt{\widehat{\bm{v}}_t}+\eps_{\mathrm{Adam}})$.
\EndFor
\end{algorithmic}
\end{algorithm}

Все операции в последней строке выполняются покомпонентно.
Поправки $1-\beta_1^t$ и $1-\beta_2^t$ устраняют начальное смещение
оценок моментов к нулю.

\subsection{Инициализация параметров}

Если все нейроны одного слоя инициализировать одинаковыми весами,
они будут получать одинаковые градиенты и останутся симметричными.
Поэтому веса инициализируют случайно.

Для активаций типа $\tanh$ часто используется инициализация Xavier:
\[
    \operatorname{Var}(w)
    \approx
    \frac{2}{n_{\mathrm{in}}+n_{\mathrm{out}}}.
\]

Для ReLU используется инициализация He:
\[
    \operatorname{Var}(w)
    \approx
    \frac{2}{n_{\mathrm{in}}}.
\]

Цель таких инициализаций --- не допустить систематического роста или
затухания дисперсии активаций и градиентов при прохождении через слои.

\subsection{Затухающие и взрывающиеся градиенты}

При обратном распространении градиент содержит произведение матриц
Якоби:
\[
    \frac{\partial J}{\partial\bm{a}^{(0)}}
    =
    \frac{\partial\bm{a}^{(1)}}{\partial\bm{a}^{(0)}}
    \cdots
    \frac{\partial\bm{a}^{(L)}}{\partial\bm{a}^{(L-1)}}
    \frac{\partial J}{\partial\bm{a}^{(L)}}.
\]

Если нормы большинства сомножителей меньше единицы, возникает
\textbf{затухание градиента}. Если они больше единицы, возникает
\textbf{взрыв градиента}.

Основные способы борьбы:

\begin{itemize}
    \item ReLU и её модификации;
    \item корректная инициализация;
    \item нормализация активаций;
    \item остаточные соединения;
    \item ограничение нормы градиента:
    \[
        \bm{g}
        \gets
        \bm{g}
        \min\left(1,\frac{\tau}{\|\bm{g}\|}\right);
    \]
    \item специализированные рекуррентные ячейки LSTM и GRU.
\end{itemize}

% ============================================================
\section{Регуляризация нейронных сетей}
% ============================================================

\subsection{$L_2$-регуляризация}

К функции потерь добавляется штраф:
\[
    J_{\mathrm{reg}}
    =
    \widehat{R}(\bm{\theta})
    +
    \frac{\lambda}{2}\sum_{\ell}\|W^{(\ell)}\|_F^2.
\]
Градиент по матрице весов:
\[
    \nabla_WJ_{\mathrm{reg}}
    =
    \nabla_W\widehat{R}+\lambda W.
\]

При обычном градиентном спуске:
\[
    W
    \gets
    (1-\eta\lambda)W
    -
    \eta\nabla_W\widehat{R}.
\]
Отсюда происходит название \textit{weight decay}.

\subsection{$L_1$-регуляризация}

Штраф
\[
    \lambda\sum_j|\theta_j|
\]
способствует разреженности параметров. Производная заменяется
субградиентом:
\[
    \partial|\theta|
    =
    \begin{cases}
        \{1\},&\theta>0,\\
        [-1,1],&\theta=0,\\
        \{-1\},&\theta<0.
    \end{cases}
\]

\subsection{Dropout}

Во время обучения компоненты вектора активаций случайно зануляются:
\[
    \widetilde{\bm{a}}
    =
    \frac{\bm{m}\odot\bm{a}}{1-p},
    \qquad
    m_j\sim\operatorname{Bernoulli}(1-p),
\]
где $p$ --- вероятность отключения нейрона.

Множитель $1/(1-p)$ обеспечивает
\[
    \E[\widetilde{\bm{a}}]=\bm{a}.
\]
При применении модели dropout отключается.

\subsection{Ранняя остановка}

После каждой эпохи контролируется значение функции потерь на
валидационной выборке. Обучение прекращается, если это значение
не улучшается в течение заданного числа эпох. Сохраняются параметры,
показавшие наилучшее валидационное качество.

\subsection{Нормализация по мини-пакету}

Для активаций одного признака вычисляются
\[
    \mu_{\mathcal{B}}
    =
    \frac{1}{B}\sum_{i=1}^{B}x_i,
\]
\[
    \sigma_{\mathcal{B}}^2
    =
    \frac{1}{B}\sum_{i=1}^{B}(x_i-\mu_{\mathcal{B}})^2.
\]
Нормализованное значение:
\[
    \widehat{x}_i
    =
    \frac{x_i-\mu_{\mathcal{B}}}
         {\sqrt{\sigma_{\mathcal{B}}^2+\eps}}.
\]
Затем применяется обучаемое преобразование
\[
    y_i=\gamma\widehat{x}_i+\beta.
\]
Во время применения модели используются накопленные оценки среднего
и дисперсии.

% ============================================================
\section{Основные архитектуры нейронных сетей}
% ============================================================

\subsection{Свёрточные сети}

Свёрточные сети предназначены прежде всего для данных с локальной
пространственной структурой. Для двумерного входа:
\[
    y_{c',i,j}
    =
    \varphi\left(
        b_{c'}
        +
        \sum_{c}
        \sum_{u}
        \sum_{v}
        K_{c',c,u,v}
        x_{c,i+u,j+v}
    \right).
\]

Основные свойства свёрточного слоя:

\begin{itemize}
    \item локальное рецептивное поле;
    \item совместное использование весов;
    \item существенно меньше параметров, чем в полном связном слое;
    \item эквивариантность сдвигу;
    \item иерархическое выделение признаков.
\end{itemize}

\subsection{Рекуррентные сети}

Рекуррентная сеть обрабатывает последовательность
$\bm{x}_1,\ldots,\bm{x}_T$:
\[
    \bm{h}_t
    =
    \varphi(
        W_{xh}\bm{x}_t
        +
        W_{hh}\bm{h}_{t-1}
        +
        \bm{b}_h
    ),
\]
\[
    \bm{y}_t
    =
    W_{hy}\bm{h}_t+\bm{b}_y.
\]

Одни и те же параметры используются на всех временных шагах.
Обучение выполняется обратным распространением ошибки во времени.
Длинные произведения матриц Якоби приводят к затухающим и
взрывающимся градиентам.

\subsection{Механизм внимания}

Пусть входные представления собраны в матрицу $X$. Строятся запросы,
ключи и значения:
\[
    Q=XW_Q,
    \qquad
    K=XW_K,
    \qquad
    V=XW_V.
\]
Масштабированное скалярное внимание:
\[
    \operatorname{Attention}(Q,K,V)
    =
    \softmax\left(
        \frac{QK^{\T}}{\sqrt{d_k}}
    \right)V.
\]
Матрица softmax определяет, с какими весами каждый элемент
последовательности агрегирует значения остальных элементов.

% ============================================================
\section{Метод опорных векторов}
% ============================================================

\subsection{Геометрия линейного классификатора}

Рассмотрим гиперплоскость
\[
    H
    =
    \{\bm{x}\in\R^d:
    \bm{w}^{\T}\bm{x}+b=0\}.
\]
Вектор $\bm{w}$ перпендикулярен гиперплоскости.

Расстояние от точки $\bm{x}$ до $H$ равно
\[
    \dist(\bm{x},H)
    =
    \frac{|\bm{w}^{\T}\bm{x}+b|}
         {\|\bm{w}\|}.
\]

Действительно, ортогональная проекция точки имеет вид
\[
    \bm{x}_{H}
    =
    \bm{x}
    -
    \frac{\bm{w}^{\T}\bm{x}+b}{\|\bm{w}\|^2}\bm{w}.
\]
Тогда
\[
    \|\bm{x}-\bm{x}_{H}\|
    =
    \left\|
        \frac{\bm{w}^{\T}\bm{x}+b}{\|\bm{w}\|^2}\bm{w}
    \right\|
    =
    \frac{|\bm{w}^{\T}\bm{x}+b|}{\|\bm{w}\|}.
\]

\begin{figure}[h]
\centering
\begin{tikzpicture}[scale=1.1,>=Latex]
    \draw[->] (-3,0) -- (3,0) node[right] {$x_1$};
    \draw[->] (0,-2) -- (0,2) node[above] {$x_2$};

    \draw[thick] (0,-1.8) -- (0,1.8)
        node[above right] {$\bm{w}^{\T}\bm{x}+b=0$};
    \draw[dashed] (-1,-1.8) -- (-1,1.8)
        node[above left] {$=-1$};
    \draw[dashed] (1,-1.8) -- (1,1.8)
        node[above right] {$=1$};

    \foreach \x/\y in {-2.4/1.1,-1.8/0.3,-2.2/-0.8,-1.2/-1.3}
    {
        \draw[red,thick] (\x-0.08,\y-0.08) -- (\x+0.08,\y+0.08);
        \draw[red,thick] (\x-0.08,\y+0.08) -- (\x+0.08,\y-0.08);
    }

    \foreach \x/\y in {2.3/1.2,1.5/0.4,2.2/-0.7,1.1/-1.2}
        \fill[blue] (\x,\y) circle (2.2pt);

    \draw[<->] (-1,-1.6) -- (1,-1.6)
        node[midway,below] {$2/\|\bm{w}\|$};
\end{tikzpicture}
\caption{Разделяющая гиперплоскость и полоса зазора SVM.}
\end{figure}

\subsection{Функциональный и геометрический отступы}

Для объекта $(\bm{x}_i,y_i)$ функциональный отступ:
\[
    \widehat{\gamma}_i
    =
    y_i(\bm{w}^{\T}\bm{x}_i+b).
\]
Он положителен при правильной классификации.

Геометрический отступ:
\[
    \gamma_i
    =
    \frac{y_i(\bm{w}^{\T}\bm{x}_i+b)}
         {\|\bm{w}\|}.
\]

Функциональный отступ зависит от масштаба параметров:
\[
    (\bm{w},b)\mapsto(c\bm{w},cb)
\]
умножает его на $c>0$, не изменяя гиперплоскость. Геометрический
отступ от такого масштабирования не зависит.

Отступ всей выборки:
\[
    \gamma
    =
    \min_i
    \frac{y_i(\bm{w}^{\T}\bm{x}_i+b)}
         {\|\bm{w}\|}.
\]

\subsection{SVM с жёстким зазором}

Для линейно разделимой выборки можно выбрать масштаб параметров так,
чтобы
\[
    y_i(\bm{w}^{\T}\bm{x}_i+b)\geq1.
\]
Тогда задача максимизации геометрического зазора принимает вид
\[
\boxed{
\begin{aligned}
    &\min_{\bm{w},b}
    &&\frac{1}{2}\|\bm{w}\|^2,\\
    &\text{при условиях}
    &&y_i(\bm{w}^{\T}\bm{x}_i+b)\geq1,
    \quad i=1,\ldots,n.
\end{aligned}
}
\]

\begin{theorem}[Эквивалентность SVM и максимизации зазора]
Пусть выборка линейно разделима и содержит объекты обоих классов.
Тогда задача SVM с жёстким зазором имеет решение
$(\bm{w}^{*},b^{*})$. Вектор $\bm{w}^{*}$ единственен, и соответствующая
гиперплоскость максимизирует геометрический отступ
\[
    \gamma
    =
    \min_i
    \frac{y_i(\bm{w}^{\T}\bm{x}_i+b)}
         {\|\bm{w}\|}.
\]
Максимальный отступ равен
\[
    \gamma^{*}=\frac{1}{\|\bm{w}^{*}\|},
\]
а ширина полосы между гиперплоскостями
$\bm{w}^{\T}\bm{x}+b=\pm1$ равна
\[
    \frac{2}{\|\bm{w}^{*}\|}.
\]
\end{theorem}

\begin{proof}
Сначала докажем существование решения. По линейной разделимости
существуют $\bm{w}_0,b_0$ такие, что
\[
    y_i(\bm{w}_0^{\T}\bm{x}_i+b_0)>0
\]
для всех $i$. Выборка конечна, поэтому
\[
    m
    =
    \min_i
    y_i(\bm{w}_0^{\T}\bm{x}_i+b_0)>0.
\]
После масштабирования
\[
    \bm{w}=\frac{\bm{w}_0}{m},
    \qquad
    b=\frac{b_0}{m}
\]
получаем допустимую точку:
\[
    y_i(\bm{w}^{\T}\bm{x}_i+b)\geq1.
\]
Следовательно, допустимое множество непусто.

Возьмём минимизирующую последовательность
$(\bm{w}_k,b_k)$. Значения $\|\bm{w}_k\|$ ограничены, поскольку
значения целевой функции сходятся к конечному инфимуму.

Выберем один положительный объект $\bm{x}_+$ и один отрицательный
объект $\bm{x}_-$. Из ограничений следует
\[
    \bm{w}_k^{\T}\bm{x}_+ + b_k\geq1,
\]
\[
    \bm{w}_k^{\T}\bm{x}_- + b_k\leq-1.
\]
Значит,
\[
    1-\bm{w}_k^{\T}\bm{x}_+
    \leq
    b_k
    \leq
    -1-\bm{w}_k^{\T}\bm{x}_-.
\]
Поскольку $\bm{w}_k$ ограничены, последовательность $b_k$ также
ограничена. Поэтому существует сходящаяся подпоследовательность
\[
    (\bm{w}_{k_j},b_{k_j})
    \to
    (\bm{w}^{*},b^{*}).
\]
Допустимое множество замкнуто, поэтому предел допустим.
Непрерывность нормы показывает, что предел достигает минимума.

Докажем единственность $\bm{w}^{*}$. Предположим, что существуют два
решения $(\bm{w}_1,b_1)$ и $(\bm{w}_2,b_2)$ с
$\bm{w}_1\neq\bm{w}_2$. Их среднее также допустимо вследствие
линейности ограничений:
\[
    \left(
        \frac{\bm{w}_1+\bm{w}_2}{2},
        \frac{b_1+b_2}{2}
    \right).
\]
Строгая выпуклость функции $\|\bm{w}\|^2$ даёт
\[
    \left\|
        \frac{\bm{w}_1+\bm{w}_2}{2}
    \right\|^2
    <
    \frac{\|\bm{w}_1\|^2+\|\bm{w}_2\|^2}{2},
\]
что противоречит оптимальности обоих решений. Следовательно,
$\bm{w}^{*}$ единственен. Смещение $b^{*}$ в вырожденных случаях
может быть неединственным.

Теперь рассмотрим произвольную разделяющую гиперплоскость
$(\bm{w},b)$. Обозначим её минимальный функциональный отступ:
\[
    m
    =
    \min_i y_i(\bm{w}^{\T}\bm{x}_i+b)>0.
\]
Её геометрический отступ равен
\[
    \gamma=\frac{m}{\|\bm{w}\|}.
\]
Масштабируем параметры:
\[
    \widetilde{\bm{w}}=\frac{\bm{w}}{m},
    \qquad
    \widetilde{b}=\frac{b}{m}.
\]
Тогда
\[
    y_i(
        \widetilde{\bm{w}}^{\T}\bm{x}_i
        +
        \widetilde{b}
    )
    \geq1,
\]
а
\[
    \|\widetilde{\bm{w}}\|
    =
    \frac{\|\bm{w}\|}{m}
    =
    \frac{1}{\gamma}.
\]
Следовательно, максимизация $\gamma$ эквивалентна минимизации
$\|\widetilde{\bm{w}}\|$ при канонических ограничениях.

В оптимальном решении минимальный функциональный отступ равен единице.
Если бы он был строго больше единицы, параметры можно было бы
уменьшить в одинаковое число раз, сохранив допустимость и уменьшив
норму $\bm{w}$. Поэтому
\[
    \gamma^{*}
    =
    \frac{1}{\|\bm{w}^{*}\|}.
\]

Расстояние между гиперплоскостями
\[
    \bm{w}^{\T}\bm{x}+b=1
    \quad\text{и}\quad
    \bm{w}^{\T}\bm{x}+b=-1
\]
равно сумме их расстояний до центральной гиперплоскости:
\[
    \frac{1}{\|\bm{w}\|}
    +
    \frac{1}{\|\bm{w}\|}
    =
    \frac{2}{\|\bm{w}\|}.
\]
\end{proof}

\subsection{Двойственная задача}

Лагранжиан задачи с жёстким зазором:
\[
    L(\bm{w},b,\bm{\alpha})
    =
    \frac{1}{2}\|\bm{w}\|^2
    -
    \sum_{i=1}^{n}
    \alpha_i
    \left[
        y_i(\bm{w}^{\T}\bm{x}_i+b)-1
    \right],
\]
где $\alpha_i\geq0$.

Условие стационарности по $\bm{w}$:
\[
    \frac{\partial L}{\partial\bm{w}}
    =
    \bm{w}
    -
    \sum_{i=1}^{n}\alpha_i y_i\bm{x}_i
    =
    \bm{0}.
\]
Следовательно,
\[
    \boxed{
    \bm{w}
    =
    \sum_{i=1}^{n}\alpha_i y_i\bm{x}_i.
    }
\]

Условие стационарности по $b$:
\[
    \frac{\partial L}{\partial b}
    =
    -
    \sum_{i=1}^{n}\alpha_i y_i
    =
    0.
\]
То есть
\[
    \boxed{
    \sum_{i=1}^{n}\alpha_i y_i=0.
    }
\]

После подстановки получаем двойственную задачу:
\[
\boxed{
\begin{aligned}
    &\max_{\bm{\alpha}}
    &&
    \sum_{i=1}^{n}\alpha_i
    -
    \frac{1}{2}
    \sum_{i=1}^{n}
    \sum_{j=1}^{n}
    \alpha_i\alpha_jy_iy_j
    \bm{x}_i^{\T}\bm{x}_j,\\
    &\text{при условиях}
    &&
    \alpha_i\geq0,\\
    &&
    &\sum_{i=1}^{n}\alpha_i y_i=0.
\end{aligned}
}
\]

Условия Каруша--Куна--Таккера включают:

\[
    y_i(\bm{w}^{\T}\bm{x}_i+b)-1\geq0,
\]
\[
    \alpha_i\geq0,
\]
\[
    \alpha_i
    \left[
        y_i(\bm{w}^{\T}\bm{x}_i+b)-1
    \right]
    =0.
\]

Последнее условие называется дополняющей нежёсткостью. Если
$\alpha_i>0$, то
\[
    y_i(\bm{w}^{\T}\bm{x}_i+b)=1.
\]
Такие объекты называются \textbf{опорными векторами}.

Решающее правило:
\[
    f(\bm{x})
    =
    \sum_{i=1}^{n}
    \alpha_i y_i
    \bm{x}_i^{\T}\bm{x}
    +
    b,
\]
\[
    \widehat{y}=\sign f(\bm{x}).
\]
Объекты с $\alpha_i=0$ не входят в решающую функцию.

\subsection{SVM с мягким зазором}

Для линейно неразделимой выборки вводятся переменные послабления
$\xi_i\geq0$:
\[
    y_i(\bm{w}^{\T}\bm{x}_i+b)
    \geq
    1-\xi_i.
\]

Задача имеет вид
\[
\boxed{
\begin{aligned}
    &\min_{\bm{w},b,\bm{\xi}}
    &&
    \frac{1}{2}\|\bm{w}\|^2
    +
    C\sum_{i=1}^{n}\xi_i,\\
    &\text{при условиях}
    &&
    y_i(\bm{w}^{\T}\bm{x}_i+b)
    \geq1-\xi_i,\\
    &&
    &\xi_i\geq0.
\end{aligned}
}
\]

Интерпретация $\xi_i$:

\begin{itemize}
    \item $\xi_i=0$ --- объект лежит на границе зазора или вне полосы;
    \item $0<\xi_i<1$ --- объект правильно классифицирован, но находится
    внутри полосы;
    \item $\xi_i=1$ --- объект находится на разделяющей гиперплоскости;
    \item $\xi_i>1$ --- объект классифицирован ошибочно.
\end{itemize}

Параметр $C>0$ определяет компромисс:

\begin{itemize}
    \item большое $C$ сильнее штрафует нарушения;
    \item малое $C$ допускает больше нарушений ради более широкого
    зазора и более сильной регуляризации.
\end{itemize}

\begin{proposition}[Эквивалентность мягкого SVM и hinge loss]
Задача SVM с мягким зазором эквивалентна безусловной задаче
\[
    \min_{\bm{w},b}
    \left[
        \frac{1}{2}\|\bm{w}\|^2
        +
        C\sum_{i=1}^{n}
        \max\left(
            0,
            1-y_i(\bm{w}^{\T}\bm{x}_i+b)
        \right)
    \right].
\]
\end{proposition}

\begin{proof}
Зафиксируем $\bm{w}$ и $b$. Для каждого $i$ требуется выбрать
минимальное $\xi_i$, удовлетворяющее условиям
\[
    \xi_i\geq0
\]
и
\[
    \xi_i
    \geq
    1-y_i(\bm{w}^{\T}\bm{x}_i+b).
\]
Так как коэффициент $C$ положителен, оптимальное значение $\xi_i$
является наименьшим допустимым:
\[
    \xi_i^{*}
    =
    \max\left(
        0,
        1-y_i(\bm{w}^{\T}\bm{x}_i+b)
    \right).
\]
Подстановка этих значений в целевую функцию даёт требуемую
безусловную задачу.
\end{proof}

Функция
\[
    \ell_{\mathrm{hinge}}(y,f)
    =
    \max(0,1-yf)
\]
называется \textbf{hinge loss}.

Двойственная задача для мягкого SVM:
\[
\boxed{
\begin{aligned}
    &\max_{\bm{\alpha}}
    &&
    \sum_{i=1}^{n}\alpha_i
    -
    \frac{1}{2}
    \sum_{i,j=1}^{n}
    \alpha_i\alpha_jy_iy_j
    \bm{x}_i^{\T}\bm{x}_j,\\
    &\text{при условиях}
    &&
    0\leq\alpha_i\leq C,\\
    &&
    &\sum_{i=1}^{n}\alpha_i y_i=0.
\end{aligned}
}
\]

Для оптимального решения выполняется:

\[
\begin{array}{lll}
    \alpha_i=0
    &\Longrightarrow&
    y_if(\bm{x}_i)\geq1,\\[2mm]
    0<\alpha_i<C
    &\Longrightarrow&
    y_if(\bm{x}_i)=1,\\[2mm]
    \alpha_i=C
    &\Longrightarrow&
    y_if(\bm{x}_i)\leq1.
\end{array}
\]

Если $0<\alpha_i<C$, то смещение можно найти по формуле
\[
    b
    =
    y_i
    -
    \sum_{j=1}^{n}
    \alpha_jy_j\bm{x}_j^{\T}\bm{x}_i.
\]
Для численной устойчивости значения $b$, полученные по нескольким
свободным опорным векторам, усредняют.

\subsection{Минимальный пример линейного SVM}

Рассмотрим одномерную выборку
\[
    x_1=-1,\quad y_1=-1,
    \qquad
    x_2=1,\quad y_2=1.
\]
Прямая задача:
\[
    \min_{w,b}\frac{1}{2}w^2
\]
при ограничениях
\[
    -1(-w+b)\geq1,
    \qquad
    w+b\geq1.
\]
Первое ограничение:
\[
    w-b\geq1.
\]
Складывая ограничения, получаем
\[
    2w\geq2,
    \qquad
    w\geq1.
\]
Минимум достигается при
\[
    w^{*}=1,
    \qquad
    b^{*}=0.
\]
Классификатор:
\[
    \widehat{y}=\sign x.
\]
Геометрический зазор:
\[
    \gamma=\frac{1}{|w^{*}|}=1.
\]

Для двойственных переменных условие
\[
    -\alpha_1+\alpha_2=0
\]
даёт
\[
    \alpha_1=\alpha_2=\alpha.
\]
Кроме того,
\[
    w
    =
    \alpha_1y_1x_1+\alpha_2y_2x_2
    =
    2\alpha.
\]
Следовательно,
\[
    \alpha_1=\alpha_2=\frac{1}{2}.
\]
Оба обучающих объекта являются опорными векторами.

% ============================================================
\section{Ядерный метод}
% ============================================================

\subsection{Переход в пространство признаков}

Линейный классификатор можно строить не в исходном пространстве, а в
пространстве признаков:
\[
    \bm{x}\mapsto\phi(\bm{x})\in\mathcal{H}.
\]
Тогда
\[
    f(\bm{x})
    =
    \langle\bm{w},\phi(\bm{x})\rangle_{\mathcal{H}}+b.
\]

В двойственной задаче объекты входят только через скалярные
произведения
\[
    \langle\phi(\bm{x}_i),\phi(\bm{x}_j)\rangle_{\mathcal{H}}.
\]
Если эти произведения можно вычислять непосредственно функцией
\[
    K(\bm{x},\bm{z})
    =
    \langle\phi(\bm{x}),\phi(\bm{z})\rangle_{\mathcal{H}},
\]
явное построение $\phi(\bm{x})$ не требуется. Это называется
\textbf{ядерным трюком}.

\subsection{Критерий допустимого ядра}

\begin{theorem}[Представление положительно полуопределённого ядра]
Пусть $\mathcal{X}$ --- непустое множество, а
$K:\mathcal{X}\times\mathcal{X}\to\R$ --- симметричная функция.
Следующие условия эквивалентны.

\begin{enumerate}
    \item Для любого $n$, любых
    $x_1,\ldots,x_n\in\mathcal{X}$ и любых
    $c_1,\ldots,c_n\in\R$ выполняется
    \[
        \sum_{i=1}^{n}
        \sum_{j=1}^{n}
        c_ic_jK(x_i,x_j)\geq0.
    \]
    \item Существуют вещественное гильбертово пространство
    $\mathcal{H}$ и отображение
    $\phi:\mathcal{X}\to\mathcal{H}$ такие, что
    \[
        K(x,z)
        =
        \langle\phi(x),\phi(z)\rangle_{\mathcal{H}}.
    \]
\end{enumerate}
\end{theorem}

\begin{proof}
Сначала докажем, что второе условие влечёт первое. Пусть
\[
    K(x,z)
    =
    \langle\phi(x),\phi(z)\rangle_{\mathcal{H}}.
\]
Тогда
\begin{align*}
    \sum_{i,j=1}^{n}c_ic_jK(x_i,x_j)
    &=
    \sum_{i,j=1}^{n}
    c_ic_j
    \langle\phi(x_i),\phi(x_j)\rangle_{\mathcal{H}}\\
    &=
    \left\langle
        \sum_{i=1}^{n}c_i\phi(x_i),
        \sum_{j=1}^{n}c_j\phi(x_j)
    \right\rangle_{\mathcal{H}}\\
    &=
    \left\|
        \sum_{i=1}^{n}c_i\phi(x_i)
    \right\|_{\mathcal{H}}^2
    \geq0.
\end{align*}

Докажем обратное утверждение. Рассмотрим векторное пространство
$\mathcal{V}_0$ формальных конечных линейных комбинаций символов
$e_x$, где $x\in\mathcal{X}$:
\[
    u=\sum_{i=1}^{n}c_ie_{x_i}.
\]
На $\mathcal{V}_0$ зададим симметричную билинейную форму:
\[
    \left\langle
        \sum_{i=1}^{n}c_ie_{x_i},
        \sum_{j=1}^{m}d_je_{z_j}
    \right\rangle_0
    =
    \sum_{i=1}^{n}
    \sum_{j=1}^{m}
    c_id_jK(x_i,z_j).
\]
По первому условию
\[
    \langle u,u\rangle_0\geq0.
\]
Однако возможно, что ненулевой формальный вектор имеет нулевую
квадратичную форму. Обозначим
\[
    \mathcal{N}
    =
    \{u\in\mathcal{V}_0:
      \langle u,u\rangle_0=0\}.
\]

Покажем, что любой $u\in\mathcal{N}$ ортогонален всем
$v\in\mathcal{V}_0$. Для любого $t\in\R$:
\[
    0
    \leq
    \langle u+tv,u+tv\rangle_0
    =
    2t\langle u,v\rangle_0
    +
    t^2\langle v,v\rangle_0.
\]
Если бы $\langle u,v\rangle_0\neq0$, то при достаточно малом $t$
противоположного знака линейный член доминировал бы над квадратичным,
и правая часть стала бы отрицательной. Следовательно,
\[
    \langle u,v\rangle_0=0.
\]
Отсюда также следует, что $\mathcal{N}$ является подпространством.

Рассмотрим факторпространство
\[
    \mathcal{V}
    =
    \mathcal{V}_0/\mathcal{N}.
\]
На нём корректно определено скалярное произведение
\[
    \langle[u],[v]\rangle
    =
    \langle u,v\rangle_0.
\]
Корректность следует из того, что элементы $\mathcal{N}$
ортогональны всем векторам.

Если класс $[u]$ имеет нулевую норму, то
\[
    \langle u,u\rangle_0=0,
\]
то есть $u\in\mathcal{N}$ и $[u]=[0]$. Поэтому полученная форма является
положительно определённым скалярным произведением.

Если $\mathcal{V}$ не полно, возьмём его пополнение $\mathcal{H}$.
Определим
\[
    \phi(x)=[e_x].
\]
Тогда
\[
    \langle\phi(x),\phi(z)\rangle_{\mathcal{H}}
    =
    \langle[e_x],[e_z]\rangle
    =
    K(x,z).
\]
Требуемое гильбертово пространство построено.
\end{proof}

\subsection{Примеры ядер}

\paragraph{Линейное ядро}
\[
    K(\bm{x},\bm{z})
    =
    \bm{x}^{\T}\bm{z}.
\]

\paragraph{Полиномиальное ядро}
\[
    K(\bm{x},\bm{z})
    =
    (\gamma\bm{x}^{\T}\bm{z}+r)^d,
\]
где $\gamma\geq0$, $r\geq0$, а $d$ --- натуральное число.

\paragraph{Гауссово RBF-ядро}
\[
    K(\bm{x},\bm{z})
    =
    \exp\bigl(-\gamma\|\bm{x}-\bm{z}\|^2\bigr),
    \qquad
    \gamma>0.
\]
Его можно разложить:
\begin{align*}
    K(\bm{x},\bm{z})
    &=
    e^{-\gamma\|\bm{x}\|^2}
    e^{-\gamma\|\bm{z}\|^2}
    e^{2\gamma\bm{x}^{\T}\bm{z}}\\
    &=
    e^{-\gamma\|\bm{x}\|^2}
    e^{-\gamma\|\bm{z}\|^2}
    \sum_{m=0}^{\infty}
    \frac{(2\gamma)^m}{m!}
    (\bm{x}^{\T}\bm{z})^m.
\end{align*}
Каждое ядро $(\bm{x}^{\T}\bm{z})^m$ соответствует скалярному
произведению тензорных степеней признаков, а коэффициенты ряда
неотрицательны. Поэтому RBF-функция является допустимым ядром.

\paragraph{Сигмоидальное ядро}
\[
    K(\bm{x},\bm{z})
    =
    \tanh(\gamma\bm{x}^{\T}\bm{z}+r).
\]
Оно является положительно полуопределённым не для всех значений
параметров, поэтому требует дополнительной проверки.

\subsection{Ядерный SVM}

В двойственной задаче скалярные произведения заменяются ядром:
\[
\boxed{
\begin{aligned}
    &\max_{\bm{\alpha}}
    &&
    \sum_{i=1}^{n}\alpha_i
    -
    \frac{1}{2}
    \sum_{i,j=1}^{n}
    \alpha_i\alpha_jy_iy_j
    K(\bm{x}_i,\bm{x}_j),\\
    &\text{при условиях}
    &&
    0\leq\alpha_i\leq C,\\
    &&
    &\sum_{i=1}^{n}\alpha_i y_i=0.
\end{aligned}
}
\]

Решающая функция:
\[
    f(\bm{x})
    =
    \sum_{i=1}^{n}
    \alpha_i y_iK(\bm{x}_i,\bm{x})
    +
    b.
\]
Фактически суммирование ведётся только по опорным векторам:
\[
    f(\bm{x})
    =
    \sum_{i:\alpha_i>0}
    \alpha_i y_iK(\bm{x}_i,\bm{x})
    +
    b.
\]

% ============================================================
\section{Последовательная минимальная оптимизация}
% ============================================================

\subsection{Основная идея SMO}

Двойственная задача SVM содержит ограничение
\[
    \sum_{i=1}^{n}\alpha_i y_i=0.
\]
Поэтому нельзя произвольно изменить только одну переменную
$\alpha_i$: необходимо одновременно изменять как минимум две
переменные.

Алгоритм SMO на каждом шаге выбирает пару $(\alpha_i,\alpha_j)$,
фиксирует остальные переменные и аналитически решает двумерную
подзадачу.

Пусть
\[
    E_i=f(\bm{x}_i)-y_i
\]
--- ошибка предсказания.

\subsection{Шаги одной итерации SMO}

\begin{enumerate}[label=\arabic*.]
    \item Выбрать индекс $i$, для которого нарушены условия
    Каруша--Куна--Таккера.

    \item Выбрать второй индекс $j\neq i$. Распространённая эвристика:
    максимизировать
    \[
        |E_i-E_j|.
    \]

    \item Вычислить допустимые границы нового значения $\alpha_j$.

    Если $y_i\neq y_j$, то
    \[
        L=\max(0,\alpha_j-\alpha_i),
    \]
    \[
        H=\min(C,C+\alpha_j-\alpha_i).
    \]

    Если $y_i=y_j$, то
    \[
        L=\max(0,\alpha_i+\alpha_j-C),
    \]
    \[
        H=\min(C,\alpha_i+\alpha_j).
    \]

    \item Вычислить
    \[
        \eta
        =
        K(\bm{x}_i,\bm{x}_i)
        +
        K(\bm{x}_j,\bm{x}_j)
        -
        2K(\bm{x}_i,\bm{x}_j).
    \]

    \item Если $\eta>0$, выполнить аналитическое обновление:
    \[
        \alpha_j^{\mathrm{unc}}
        =
        \alpha_j
        +
        \frac{y_j(E_i-E_j)}{\eta}.
    \]

    \item Ограничить новое значение интервалом:
    \[
        \alpha_j^{\mathrm{new}}
        =
        \min\left(
            H,
            \max(L,\alpha_j^{\mathrm{unc}})
        \right).
    \]

    \item Восстановить $\alpha_i$ из линейного ограничения:
    \[
        \alpha_i^{\mathrm{new}}
        =
        \alpha_i
        +
        y_iy_j
        \left(
            \alpha_j-\alpha_j^{\mathrm{new}}
        \right).
    \]

    \item Обозначим
    \[
        \Delta\alpha_i
        =
        \alpha_i^{\mathrm{new}}-\alpha_i,
        \qquad
        \Delta\alpha_j
        =
        \alpha_j^{\mathrm{new}}-\alpha_j.
    \]
    Вычислить кандидаты для смещения:
    \[
    \begin{aligned}
        b_1
        &=
        b-E_i
        -y_i\Delta\alpha_iK(\bm{x}_i,\bm{x}_i)
        -y_j\Delta\alpha_jK(\bm{x}_i,\bm{x}_j),\\
        b_2
        &=
        b-E_j
        -y_i\Delta\alpha_iK(\bm{x}_i,\bm{x}_j)
        -y_j\Delta\alpha_jK(\bm{x}_j,\bm{x}_j).
    \end{aligned}
    \]

    \item Обновить $b$:
    \[
        b=
        \begin{cases}
            b_1,
            &0<\alpha_i^{\mathrm{new}}<C,\\
            b_2,
            &0<\alpha_j^{\mathrm{new}}<C,\\
            (b_1+b_2)/2,
            &\text{иначе}.
        \end{cases}
    \]

    \item Повторять итерации, пока нарушения условий
    Каруша--Куна--Таккера не станут меньше заданного допуска.
\end{enumerate}

\subsection{Минимальный пример шага SMO}

Возьмём выборку
\[
    x_1=-1,\quad y_1=-1,
    \qquad
    x_2=1,\quad y_2=1
\]
и линейное ядро $K(x,z)=xz$. Пусть
\[
    \alpha_1=\alpha_2=0,
    \qquad
    b=0,
    \qquad
    C\geq\frac{1}{2}.
\]

В начале $f(x)=0$, поэтому
\[
    E_1=f(x_1)-y_1=1,
    \qquad
    E_2=f(x_2)-y_2=-1.
\]
Так как $y_1\neq y_2$:
\[
    L=0,
    \qquad
    H=C.
\]
Кроме того,
\[
    \eta
    =
    K(-1,-1)+K(1,1)-2K(-1,1)
    =
    1+1-2(-1)=4.
\]
Новое значение:
\[
    \alpha_2^{\mathrm{new}}
    =
    0+\frac{1\cdot(1-(-1))}{4}
    =
    \frac{1}{2}.
\]
Далее
\[
    \alpha_1^{\mathrm{new}}
    =
    0+(-1)(1)\left(0-\frac{1}{2}\right)
    =
    \frac{1}{2}.
\]
Вычисление смещения даёт $b=0$. Таким образом, одна парная
оптимизация сразу находит решение
\[
    \alpha_1=\alpha_2=\frac{1}{2},
    \qquad
    f(x)=x.
\]

% ============================================================
\section{Многоклассовый SVM и регрессия}
% ============================================================

\subsection{Многоклассовая классификация}

Базовый SVM является бинарным классификатором. Для $K$ классов
применяются схемы декомпозиции.

\paragraph{Один против остальных}

Строится $K$ классификаторов. Классификатор $k$ отделяет класс $k$
от всех остальных. Предсказание:
\[
    \widehat{y}
    =
    \argmax_{k=1,\ldots,K}f_k(\bm{x}).
\]

Преимущество --- требуется только $K$ моделей. Недостаток ---
возможный дисбаланс между одним положительным классом и объединением
остальных классов.

\paragraph{Один против одного}

Для каждой пары классов строится отдельный классификатор. Всего:
\[
    \frac{K(K-1)}{2}
\]
моделей. Итоговый класс выбирается голосованием или агрегированием
значений решающих функций.

\subsection{Метод опорных векторов для регрессии}

В $\eps$-SVR используется нечувствительная к малым ошибкам функция
потерь:
\[
    \ell_{\eps}(y,f)
    =
    \max(0,|y-f|-\eps).
\]

Прямая задача:
\[
\begin{aligned}
    \min_{\bm{w},b,\bm{\xi},\bm{\xi}^{*}}
    \quad&
    \frac{1}{2}\|\bm{w}\|^2
    +
    C\sum_{i=1}^{n}(\xi_i+\xi_i^{*}),\\
    \text{при условиях}\quad&
    y_i-f(\bm{x}_i)\leq\eps+\xi_i,\\
    &
    f(\bm{x}_i)-y_i\leq\eps+\xi_i^{*},\\
    &
    \xi_i,\xi_i^{*}\geq0.
\end{aligned}
\]

Двойственная задача:
\[
\begin{aligned}
    \max_{\bm{\alpha},\bm{\alpha}^{*}}
    \quad&
    -\frac{1}{2}
    \sum_{i,j=1}^{n}
    (\alpha_i-\alpha_i^{*})
    (\alpha_j-\alpha_j^{*})
    K(\bm{x}_i,\bm{x}_j)\\
    &-
    \eps\sum_{i=1}^{n}
    (\alpha_i+\alpha_i^{*})
    +
    \sum_{i=1}^{n}
    y_i(\alpha_i-\alpha_i^{*}),\\
    \text{при условиях}\quad&
    \sum_{i=1}^{n}(\alpha_i-\alpha_i^{*})=0,\\
    &
    0\leq\alpha_i,\alpha_i^{*}\leq C.
\end{aligned}
\]

Регрессионная функция:
\[
    f(\bm{x})
    =
    \sum_{i=1}^{n}
    (\alpha_i-\alpha_i^{*})
    K(\bm{x}_i,\bm{x})
    +
    b.
\]
Объекты, для которых
$\alpha_i-\alpha_i^{*}\neq0$, являются опорными.

% ============================================================
\section{Практическое применение SVM}
% ============================================================

\subsection{Предобработка данных}

SVM чувствителен к масштабу признаков, поскольку скалярные произведения
и расстояния непосредственно входят в модель. Обычно выполняется
стандартизация:
\[
    x_j^{\mathrm{scaled}}
    =
    \frac{x_j-\mu_j}{\sigma_j}.
\]

Параметры $\mu_j$ и $\sigma_j$ вычисляются только по обучающей части.
Та же трансформация применяется к валидационным и тестовым объектам.

\subsection{Выбор гиперпараметров}

Для линейного SVM главным гиперпараметром является $C$.

Для RBF-SVM выбираются:

\[
    C>0,
    \qquad
    \gamma>0.
\]

Параметр $\gamma$ определяет ширину ядра:

\begin{itemize}
    \item малое $\gamma$ даёт гладкую границу;
    \item большое $\gamma$ делает влияние каждого объекта локальным
    и повышает риск переобучения.
\end{itemize}

Параметры обычно выбираются кросс-валидацией по логарифмической сетке,
например среди значений
\[
    C\in\{10^{-3},10^{-2},\ldots,10^3\},
\]
\[
    \gamma\in\{10^{-4},10^{-3},\ldots,10^2\}.
\]

\subsection{Несбалансированные классы}

Если классы имеют разные размеры или разную цену ошибки, используют
разные штрафы:
\[
    C_+
    \quad\text{и}\quad
    C_-.
\]
Прямая задача:
\[
    \min
    \frac{1}{2}\|\bm{w}\|^2
    +
    C_+\sum_{i:y_i=+1}\xi_i
    +
    C_-\sum_{i:y_i=-1}\xi_i.
\]

\subsection{Шаги построения SVM}

\begin{enumerate}
    \item Разделить данные на обучающую, валидационную и тестовую
    части.
    \item Обработать пропущенные значения.
    \item Масштабировать признаки.
    \item Выбрать линейное или нелинейное ядро.
    \item Выбрать сетку гиперпараметров.
    \item Выполнить кросс-валидацию только на обучающих данных.
    \item Обучить модель с лучшими параметрами на полной обучающей
    части.
    \item Один раз оценить качество на тестовой выборке.
    \item Проанализировать число опорных векторов и ошибки по классам.
\end{enumerate}

% ============================================================
\section{Сопоставление нейронных сетей и SVM}
% ============================================================

\begin{center}
\begin{longtable}{p{3.4cm} p{5.6cm} p{5.6cm}}
\toprule
\textbf{Критерий}
&
\textbf{Нейронные сети}
&
\textbf{SVM}
\\
\midrule

Целевая функция
&
Обычно невыпуклая по всем параметрам.
&
Выпуклая задача квадратичного программирования.
\\
\midrule

Оптимум
&
Градиентный алгоритм не обязан находить глобальный минимум.
&
При корректной постановке находится глобальное решение; вектор
нормали линейного SVM единственен.
\\
\midrule

Представление признаков
&
Признаки обучаются совместно с классификатором.
&
Обычно признаки задаются заранее или неявно определяются ядром.
\\
\midrule

Объём данных
&
Особенно эффективны при больших объёмах данных.
&
Часто эффективны на малых и средних выборках.
\\
\midrule

Высокая размерность
&
Требуется контроль числа параметров и регуляризация.
&
Линейный SVM хорошо работает с высокоразмерными разреженными
признаками.
\\
\midrule

Нелинейность
&
Создаётся композицией слоёв и функций активации.
&
Создаётся ядром.
\\
\midrule

Стоимость предсказания
&
Определяется архитектурой и числом параметров.
&
Для ядерного SVM пропорциональна числу опорных векторов.
\\
\midrule

Обучение
&
Хорошо масштабируется за счёт мини-пакетов и аппаратных ускорителей.
&
Полный ядерный SVM может требовать хранения матрицы Грама размера
$n\times n$.
\\
\midrule

Интерпретация регуляризации
&
Weight decay, dropout, ранняя остановка, аугментация данных.
&
Ширина зазора, параметр $C$, параметры ядра.
\\
\bottomrule
\end{longtable}
\end{center}

\subsection{Общие черты}

Оба подхода решают задачу приближения неизвестной зависимости
по конечной выборке и требуют контроля сложности модели.

Для линейного SVM:
\[
    \frac{1}{2}\|\bm{w}\|^2
    +
    C\sum_i\ell_{\mathrm{hinge}}(y_i,f_i).
\]
Для нейронной сети с $L_2$-регуляризацией:
\[
    \frac{1}{n}\sum_i
    \ell(y_i,f(\bm{x}_i;\bm{\theta}))
    +
    \frac{\lambda}{2}\|\bm{\theta}\|^2.
\]
В обоих случаях целевая функция объединяет ошибку на данных и штраф
за сложность параметров.

Последний слой нейронной сети также часто является линейным
классификатором:
\[
    f(\bm{x})
    =
    \bm{w}^{\T}\bm{h}(\bm{x})+b,
\]
где $\bm{h}(\bm{x})$ --- представление, построенное предыдущими слоями.
Таким образом, нейронную сеть можно рассматривать как совместное
обучение признакового отображения и классификатора.

% ============================================================
\section{Типичные ошибки и контрольные вопросы}
% ============================================================

\subsection{Типичные ошибки при работе с нейронными сетями}

\begin{itemize}
    \item Использование линейных активаций во всех слоях.
    \item Инициализация всех весов одинаковыми значениями.
    \item Применение softmax отдельно от численно устойчивой
    кросс-энтропии.
    \item Настройка модели по тестовой выборке.
    \item Отсутствие нормализации входных признаков.
    \item Слишком большая скорость обучения.
    \item Отсутствие контроля переобучения.
    \item Обновление весов во время вычисления градиентов одного
    обратного прохода вместо одновременного обновления после
    вычисления всех градиентов.
\end{itemize}

\subsection{Типичные ошибки при работе с SVM}

\begin{itemize}
    \item Отсутствие масштабирования признаков.
    \item Интерпретация параметра $C$ как коэффициента регуляризации
    в обычном смысле: увеличение $C$ ослабляет регуляризацию,
    поскольку сильнее штрафует ошибки.
    \item Использование произвольной симметричной функции вместо
    положительно полуопределённого ядра.
    \item Настройка $C$ и $\gamma$ по тестовой выборке.
    \item Предположение, что все обучающие объекты входят в итоговую
    решающую функцию.
    \item Предположение, что SVM непосредственно выдаёт
    калиброванные вероятности классов.
\end{itemize}

\subsection{Краткая сводка формул}

\paragraph{Перцептрон}
\[
    \widehat{y}
    =
    \sign(\bm{w}^{\T}\bm{x}+b),
\]
\[
    \bm{w}\gets\bm{w}+\eta y_i\bm{x}_i,
    \qquad
    b\gets b+\eta y_i
\]
при ошибке.

\paragraph{Полносвязный слой}
\[
    \bm{z}^{(\ell)}
    =
    W^{(\ell)}\bm{a}^{(\ell-1)}
    +
    \bm{b}^{(\ell)},
\]
\[
    \bm{a}^{(\ell)}
    =
    \varphi^{(\ell)}(\bm{z}^{(\ell)}).
\]

\paragraph{Обратное распространение}
\[
    \bm{\delta}^{(\ell)}
    =
    (W^{(\ell+1)})^{\T}
    \bm{\delta}^{(\ell+1)}
    \odot
    \varphi'^{(\ell)}(\bm{z}^{(\ell)}),
\]
\[
    \nabla_{W^{(\ell)}}J
    =
    \bm{\delta}^{(\ell)}
    (\bm{a}^{(\ell-1)})^{\T}.
\]

\paragraph{Жёсткий SVM}
\[
    \min_{\bm{w},b}
    \frac{1}{2}\|\bm{w}\|^2,
    \qquad
    y_i(\bm{w}^{\T}\bm{x}_i+b)\geq1.
\]

\paragraph{Мягкий SVM}
\[
    \min_{\bm{w},b}
    \frac{1}{2}\|\bm{w}\|^2
    +
    C\sum_i
    \max(0,1-y_if(\bm{x}_i)).
\]

\paragraph{Двойственная задача}
\[
    \max_{\bm{\alpha}}
    \left[
        \sum_i\alpha_i
        -
        \frac{1}{2}
        \sum_{i,j}
        \alpha_i\alpha_jy_iy_jK(\bm{x}_i,\bm{x}_j)
    \right],
\]
\[
    0\leq\alpha_i\leq C,
    \qquad
    \sum_i\alpha_i y_i=0.
\]

\paragraph{Ядерный классификатор}
\[
    f(\bm{x})
    =
    \sum_{i:\alpha_i>0}
    \alpha_i y_iK(\bm{x}_i,\bm{x})
    +
    b,
\]
\[
    \widehat{y}=\sign f(\bm{x}).
\]

\end{document}
